\clearpage
\section{Node Embedding and Its Relation to Matrix Factorization}

\subsection{Simple Matrix Factorization}

For the objective
\[
  \min_Z \norm{A-Z^{\mathsf T}Z}_2,
\]
identify the decoder in the encoder--decoder view of node embeddings.

\begin{solution}
Let \(Z=[\vect z_1,\ldots,\vect z_n]\in\R^{d\times n}\), so node embeddings
are the columns of \(Z\).  Since
\[
  (Z^{\mathsf T}Z)_{ij}=\vect z_i^{\mathsf T}\vect z_j,
\]
the decoder is the inner product
\[
  \boxed{\operatorname{DEC}(\vect z_i,\vect z_j)
  =\vect z_i^{\mathsf T}\vect z_j}.
\]
It predicts \(A_{ij}\), and the stated objective penalizes the difference
between all predicted and observed pairwise similarities.
\end{solution}

\subsection{Alternate Matrix Factorization}

Suppose the decoder is the bilinear form
$\vect{z}_i^{\mathsf T}W\vect{z}_j$. Write the corresponding matrix
factorization objective.

\begin{solution}
For \(W\in\R^{d\times d}\), the reconstructed entry and matrix are
\[
  \widehat A_{ij}=\vect z_i^{\mathsf T}W\vect z_j,
  \qquad
  \widehat A=Z^{\mathsf T}WZ.
\]
Thus, when both the embeddings and decoder are learned, the objective is
\[
  \boxed{\min_{Z,W}\norm{A-Z^{\mathsf T}WZ}_2}.
\]
If \(W\) is fixed, the minimization is over \(Z\) alone.
\end{solution}

\subsection{Relation to Eigendecomposition}

State the condition on $W$ under which the factorization in 5.2 is equivalent
to learning an eigendecomposition of $A$.

\begin{solution}
For a real symmetric adjacency matrix, write its eigendecomposition as
\[
  A=U\Lambda U^{\mathsf T},
  \qquad U^{\mathsf T}U=I,
\]
where \(\Lambda\) is diagonal and may contain negative eigenvalues.  Therefore
\(Z^{\mathsf T}WZ\) has this form when
\[
  \boxed{W\text{ is diagonal},\qquad ZZ^{\mathsf T}=I}.
\]
Specifically, \(U=Z^{\mathsf T}\) contains the eigenvectors and
\(W=\Lambda\) contains the corresponding eigenvalues.  If \(d<n\), this is
the analogous rank-\(d\) truncated eigendecomposition.
\end{solution}

\subsection{Multi-Hop Node Similarity}

Nodes are similar when connected by at least one path of length at most $k$.
Using an embedding lookup encoder and inner-product decoder, formulate the
corresponding matrix factorization problem.

\begin{solution}
Define the binary \(k\)-hop similarity matrix by
\[
  S^{(k)}_{ij}
  =
  \begin{cases}
    1, & i\ne j\text{ and }
      \left(\sum_{\ell=1}^{k}A^\ell\right)_{ij}>0,\\
    0, & \text{otherwise}.
  \end{cases}
\]
The powers of \(A\) detect walks of each length, and thresholding implements
the requirement that at least one path of length at most \(k\) exists.  With
the inner-product decoder, the desired matrix factorization is
\[
  \boxed{\min_Z\norm{S^{(k)}-Z^{\mathsf T}Z}_2}.
\]
\end{solution}

\subsection{Limitation of node2vec (i)}

Characterize the node2vec representations of the 10-node cliques in the
provided graph. Explain whether those embeddings reflect structural similarity.

\begin{solution}
Let \(C_L\) and \(C_R\) be the left and right 10-node cliques.  Because every
pair of nodes in a clique is adjacent, node2vec walks generate many
co-occurrences within the same clique.  Hence it tends to produce
\[
  \vect z_u\approx\vect z_v
  \quad\text{for }u,v\in C_L,
  \qquad\text{or for }u,v\in C_R.
\]
However, the two pictured components are disconnected, so an original-graph
random walk can never move from one clique to the other.  Consequently there
are no positive cross-component co-occurrences to align a node in \(C_L\)
with its structurally corresponding node in \(C_R\); negative sampling may
even push them apart.  Thus node2vec mainly captures proximity, and these
embeddings need not reflect the identical structural roles across the two
cliques.
\end{solution}

\subsection{Limitation of node2vec (ii)}

After arriving at node $w$, list the nodes reachable in one additional step
under node2vec and under struct2vec, assuming $g_k(u,v)>0$ for all $u,v,k$.

\begin{solution}
Node \(w\) is an ordinary node of the right 10-node clique.  Its original-graph
neighbors are exactly the other nine clique nodes, so node2vec can reach
\[
  \boxed{C_R\setminus\{w\}}
\]
in one step.  It cannot reach the tail directly or any node in the other
component in that step.

Struct2vec instead walks on its constructed clique over the entire original
vertex set \(V\).  Since \(g_k(u,v)>0\) for every pair, every such edge has
positive transition probability.  Therefore it can reach
\[
  \boxed{V\setminus\{w\}}
\]
in one step, using the usual convention that the walk has no self-loop.
\end{solution}

\subsection{Limitation of node2vec (iii)}

Explain why struct2vec must consider multiple structural similarity functions
$g_k$ during its random walk.

\begin{solution}
Each \(g_k\) compares structure at a different radius.  A small \(k\) captures
immediate properties such as degree and the degrees of nearby nodes, but two
nodes that look identical locally may differ beyond that radius.  Larger
\(k\) values reveal progressively more global roles, such as whether a node
lies near a tail, hub, or dense subgraph.  Using several \(g_k\)'s lets the
walk preserve structural similarity at multiple scales rather than making
all decisions from one possibly insufficient neighborhood radius.
\end{solution}

\subsection{Limitation of node2vec (iv)}

Characterize the representations of the two 10-node cliques after running
struct2vec on the graph.

\begin{solution}
Struct2vec is not restricted by original-graph connectivity; its transition
weights compare structural roles directly.  The two pictured components are
structurally isomorphic, so corresponding nodes across them receive high
structural-similarity weights and should have similar embeddings.

In particular, let \(b_L\) and \(b_R\) be the green clique nodes attached to
the two tails.  Then
\[
  \vect z_{b_L}\approx\vect z_{b_R}.
\]
Likewise, the nine ordinary clique nodes on either side share the same role
and should form one cross-clique embedding cluster.  The attachment nodes can
remain distinct from that ordinary-node cluster because their additional tail
edge gives them a different structural role.  Thus, unlike node2vec,
struct2vec aligns the two cliques by role even though they lie in disconnected
components.
\end{solution}
